\section{Local densities at all target valuations}\label{sec:local}
On a shell at \(m\), put \(A_1=N(z)\) and \(A_2=N(y)\).
The bump \(W\) restricts both norms to be positive. At a variable
prime \(p\ge5\) write \(g_m(p)\) for the density of the bad test
\(p\mid A_1A_2\). For \(p\nmid m\) this is the conditional mean in
\eqref{eq:angular}. At \(p\mid m\) we derive its analogue by reducing
to coprime Hecke indices.

\subsection{Coprime target and level}
At \(3\) the ordinary family requires no test, since \(m\equiv2\bmod3\)
forces both norms to be units modulo \(3\). In the mixed family
\(m\equiv1\bmod3\), and we require the prescribed component to have
\(3\)-adic valuation exactly one. Among the nine elements of \(R/3R\)
there are six units, two nonzero nonunits, and zero. The nonzero
nonunits have norm valuation one, since \(1-\omega\) has norm \(3\)
and its square is associated to \(3\). Norm one modulo \(3\)
requires one unit and one nonunit, giving \(2\cdot6\cdot3=36\)
raw quaternion residues. Exactly \(6\cdot2\) have a unit in one
prescribed component and a nonzero nonunit in the other. Passing to projective classes preserves the ratio,
so either mixed orientation has mean \(1/3\).

At an inert \(p\nmid m\), by anisotropy a bad point has one
component zero modulo \(pR\). Its image in
\(\PGL_2(\F_p)\) lies in the normalizer of the nonsplit torus
\((R/pR)^\times/\F_p^\times\). This is the union of that torus and
its translate by \(j\), with \(2(p+1)\) elements out of
\(p(p^2-1)\). The norms of the torus elements cover the squareclasses
equally, so \(g_m(p)=2/[p(p-1)]\) in either class.

At a split \(p\nmid m\), the two factors of the Eisenstein
subalgebra occupy the diagonal of \(M_2(\F_p)\). The other
Eisenstein component occupies the off-diagonal. Their norms are
the products of the respective entries, up to sign.
For a fixed nonzero determinant, the first row is uniform on
nonzero rows, each with \(p\) completions. Thus a given entry is
zero with probability \(1/(p+1)\). Two zero entries are possible
only when they form an opposite pair. For either such pair there
are \(p-1\) matrices of the given determinant, giving probability
\(1/[p(p+1)]\). Inclusion-exclusion gives
\begin{equation}\label{eq:B}
1-g_m(p)=B_p:=1-\frac4{p+1}+\frac2{p(p+1)}
             =\frac{(p-1)(p-2)}{p(p+1)}.
\end{equation}
For the exact condition \(v_p(A_i)=1\) at either one prescribed
component, exactly one of its two entries is zero modulo \(p\),
and this entry must be nonzero modulo \(p^2\). All matrix lifts of
an invertible residue are invertible, and determinant squareclasses
are determined modulo \(p\). The conditional probability of the
last lift condition is \((p-1)/p\). Therefore the mean of
\(v_p(A_i)=1\) at level \(p^2\) is
\begin{equation}\label{eq:t}
 t_p=\frac{2(p-1)^2}{p^2(p+1)},\qquad
 a_p:=\frac{t_p}{2B_p}=\frac{p-1}{p(p-2)} .
\end{equation}
These means hold on projective groups as well, since all tests
are unchanged by a unit scalar.

\subsection{Primes dividing the target}
For inert \(p\mid m\), divisibility of one norm by \(p\) forces
the other to be divisible by \(p\) as well, so both components
are in \(pR\). The bad set is empty if \(v_p(m)=1\).
Otherwise it is obtained by multiplying both coordinates by \(p\)
from the shell at \(m/p^2\). This preserves \(W\) and the other
prime tests. Hence for \(e=v_p(m)\) the scalar density is
\begin{equation}\label{eq:inert}
g_m(p)=0\ (e=1),\qquad
g_m(p)=\frac{\sigma(p^{e-2})}{\sigma(p^e)}\ (e\ge2).
\end{equation}

Now let \(p=\pi\bar\pi\) be split in \(R\), and fix \(r\) prime
to \(p\). At target \(p^e r\), let \(U_e\) be the weighted count
with no restriction at \(p\), and \(G_e\) the count with both
norms units at \(p\). Carry along \(W\) and the other prime
tests. At target \(r\), let \(L_0\) count with the additional
test
\[
 \frac12\bigl(\mathbf1_{\pi\mid z}+\mathbf1_{\bar\pi\mid z}
        +\mathbf1_{\pi\mid y}+\mathbf1_{\bar\pi\mid y}\bigr).
\]
It has mean \(2/(p+1)\) by the single-entry calculation.
If one norm has valuation \(v>0\), the \(v+1\) allocations of its
two prime factors are equally frequent. The average of the test
at that pair of norms is \(v/(v+1)\). The conditions at other
primes are unchanged on reallocating these factors. Indeed,
dividing out or inserting factors above \(p\) acts by units at
every other prime, preserving all the norm valuations and the
individual split-prime divisibilities used in the tests.

Classify a pair of norms on the shell \(p^e r\) by the minimum
\(j\) of its two \(p\)-adic valuations. If \(j<e\), both valuations
are \(j\); removing \(p^j\) from each norm gives a \(G_{e-j}\)
pair with allocation multiplier \((j+1)^2\). If \(j=e\), the
residual target is \(r\). When both residual norms are units
the multiplier is \((e+1)^2\); if one has valuation \(v>0\),
the ratio of allocation counts is
\((e+1)(v+e+1)/(v+1)=(e+1)^2-e(e+1)v/(v+1)\).
In either reduction the normalized norm coordinate stays fixed.
This proves
\begin{equation}\label{eq:recurrence}
 U_e=\sum_{j=0}^{e-1}(j+1)^2 G_{e-j}+(e+1)^2 U_0-e(e+1)L_0
 \qquad(e\ge1).
\end{equation}

The generating series for \((j+1)^2\) is
\(A(T)=(1+T)/(1-T)^3\), and that for \(e(e+1)\) is
\(C(T)=2T/(1-T)^3\). Thus
\[
 \sum_{e\ge1}G_eT^e=\frac{\sum_{e\ge0}U_eT^e}{A(T)}
                     -U_0+\frac{C(T)}{A(T)}L_0 .
\]
The coefficients of \(1/A(T)\) have absolute value at most eight,
and those of \(C(T)/A(T)\) at most two. The bad count \(U_e-G_e\)
is therefore a combination of unrestricted counts at lower
\(p\)-exponents and the \(L_0\) count at the unit target, with
coefficient absolute sum at most \(8e+2\). Applying the same
identity to the scalar prediction from \eqref{eq:angular} uses
\(U_e=\sigma(p^e)\), \(U_0=1\), and \(L_0=2/(p+1)\) after
factoring out the other local factors. Substitution gives
\(\sum_{e\ge1}G_e T^e=\frac{(p-1)^2}{p+1}\frac{T}{1-pT}\)
for those scalar values. In particular
\begin{equation}\label{eq:split}
 1-g_m(p)=\frac{(p-1)^2p^{e-1}}{(p+1)\sigma(p^e)}
                 \qquad(e=v_p(m)\ge1,\ p\text{ split}).
\end{equation}

\paragraph{Example at a split square.}\label{sec:split-example}
Take \(m=7^2r\) with \(7\nmid r\). The first two cases of
\eqref{eq:recurrence} give
\[
 U_2-G_2=4G_1+9U_0-6L_0=4U_1-7U_0+2L_0 .
\]
The term at \(7r\) is unrestricted at \(7\), and the only
remaining \(7\)-level test is \(L_0\) at \(r\). For the main
term, factor out the other primes and insert
\(U_2=57,\ U_1=8,\ U_0=1,\ L_0=1/4\).
The resulting main term for \(G_2\) is \(63/2\), giving
\(1-g_m(7)=(63/2)/57\).
Including the shell factor \(\sigma(7^2)/7^2=57/49\) gives
\(9/14\). This is \(B_7\cdot(6/5)\), since \(B_7=15/28\)
at a unit target. In the cofactor sum of Section~\ref{sec:cofactor-log}
this gives weight \(6/5\) without a common factor \(7\), and
\((6/5)/7\) with that common factor, since the residual target
\(7r\) still has positive valuation at \(7\). In general
\eqref{eq:split} gives \(E_p=(p-1)/(p-2)\) at every positive
valuation.

\subsection{The joint estimate}
Fix squarefree coprime \(u,v\) made of split primes coprime to
\(m\). Place the condition \(v_p(A_1)=1\) at \(p\mid u\) and
\(v_p(A_2)=1\) at \(p\mid v\). Give each coordinate pair weight
\(W(t)/\tau(uv)\). Here and below \(m\) is even and prime to \(3\).
Keep the odd-component condition and the prescribed condition
at \(3\). For any squarefree \(\ell\) coprime to \(6uv\), let
\(S_m(\ell)\) be the weighted count additionally bad at every
prime dividing \(\ell\), and put
\[
 X_m=I_\phi h R_o(m)\beta_m,\qquad
 \beta_m=\begin{cases}1,&m\equiv2\bmod6,\\1/3,&m\equiv4\bmod6.\end{cases}
\]
Extend \(g_m\) multiplicatively on squarefree products.

\begin{lemma}\label{lem:joint}
For \(m\le n\) as above and any fixed \(\zeta>0\),
\begin{equation}\label{eq:joint}
 S_m(\ell)=X_m\prod_{p\mid uv}(t_p/2)\,g_m(\ell)
   +O_{\phi,\zeta}\!\left((uv)^3\ell^{3/2}n^{1/2+\zeta}
                                           \sqrt{M/H}\right).
\end{equation}
\end{lemma}
\begin{proof}
First omit oddness. If the level is coprime to \(m\), apply
\eqref{eq:angular}. A factor at \(p\mid uv\) has size
\(p^3|\PGL_2(\F_p)|<p^6\) and a factor at \(p\mid\ell\)
has size less than \(p^3\). Thus the norm cost is
\(O((uv)^3\ell^{3/2})\), the fixed factor at \(3\) being bounded.

For \(p\mid(\ell,m)\), eliminate the test by inert scaling or
\eqref{eq:recurrence}. Each time, the remaining test at the
eliminated prime is present only at a reduced target coprime to it.
The allocation calculation works with any of the other tests
still imposed. In particular the test at a split prime produced
by an earlier elimination is an average of individual
divisibilities, all invariant under changes of allocation at
another prime. The coprime conditional means are independent of
the nonzero determinant class at those other primes. Extracting
split powers and inert squares does not change the target class
modulo \(3\), and the prescribed orientation is preserved. Hence
the scalar terms factor, with local factors
\eqref{eq:inert} and \eqref{eq:split}. The total coefficient variation
in the elimination is at most \(\tau(m)^4\), because
\(8e+2\le(e+1)^4\). Each resulting target divides \(m\), and the
remaining level divides the original level. The product of these
losses with the divisor factor in \eqref{eq:angular} is absorbed
by \(O_\zeta(n^\zeta)\).

Finally subtract the shell at \(m/4\) if it exists. Rescaling its
coordinates by \(2\) preserves the normalized angle, the odd-prime
tests and the orientation at \(3\). The valuations at all odd
primes agree with those for \(m\), so the scalar \(\sigma(m)\)
is replaced by \(\sigma(m)-\sigma(m/4)\). This gives
\eqref{eq:joint}.
\end{proof}
